\chapter{Development workflow}

\label{ch:workflow}

This chapter answers how a finding becomes a commit here, and why the answer is a queue of
rows rather than a conversation. The loop is: a finding, a spec that a cheaper model can
implement, a row that carries the spec and its criteria, a lane that dispatches a worker, a
verify step, the gate, a commit --- and, when any of those fail, a park that keeps the work
and moves the queue on. The parts that matter are the ones that keep the loop honest: the row
cannot drift from its spec, a park cannot wedge the queue, a resume cannot become a retry, and
a commit cannot name a tool.

\section{The loop}

\begin{center}
\resizebox{\linewidth}{!}{%
\begin{tikzpicture}[node distance=8mm and 12mm]
	\node[b] (finding) {a finding\\\scriptsize measured, with a path};
	\node[b, right=of finding] (spec) {a spec\\\scriptsize premises $\cdot$ criteria $\cdot$ files};
	\node[b, right=of spec] (row) {a row\\\scriptsize plan.tsv};
	\node[b, below=of row] (lane) {the lane\\\scriptsize dispatch};
	\node[b, left=of lane] (gate) {the gate\\\scriptsize 20 steps};
	\node[b, left=of gate] (commit) {the commit};
	\node[b, above=of lane] (review) {review\\\scriptsize verdict per row};

	\draw[f] (finding) -- (spec);
	\draw[f] (spec) -- (row);
	\draw[f] (row) -- (lane);
	\draw[f] (lane) -- (gate);
	\draw[f] (gate) -- (commit);
	\draw[sf] (lane) -- (review);
	\draw[bk] (review) |- node[p, right=2pt, midway]{re-spec} (spec);
\end{tikzpicture}}
\captionof{figure}{The red edge is what keeps the queue honest, because a parked row is
reviewed and returns as a better spec rather than as a retry.}
\end{center}

Each hop has a file behind it. The finding is a paragraph with a measured path and a command
that reproduces it. The spec is a numbered document under \texttt{specs/}. The row is a line
in \texttt{scripts/autonomy/plan.tsv}. The lane is a dispatch script. The gate is
\texttt{quality-gate.sh}. The commit is one commit with a neutral identity. Nothing in the
loop depends on a session's memory, which is what lets a cycle be killed and resumed without
losing the thread.

\section{The spec is written for a model that cannot guess}

A spec here has a fixed shape, and the shape is the point: a cheaper model implements it, so
every fact the implementer would otherwise have to guess is written down.

\begin{itemize}
	\item \textbf{Why this exists} --- the failure, measured, with the number that makes it a
	failure.
	\item \textbf{Verified premises} --- facts measured on a date, each with the path or
	command that produced it.
	\item \textbf{Requirements} --- behavioural, numbered, each checkable.
	\item \textbf{Files} --- the paths the work may touch.
	\item \textbf{Exit criteria} --- shell commands that fail today.
	\item \textbf{Guardrails} --- what must not be done, including the tempting wrong fix.
\end{itemize}

Spec 376 is a good example of the first two: it states that \texttt{pii-store} has no
retention code at all, and it proves it with a search whose count is zero. Spec 362 does the
same for the model ids: seven graphs hardcode a model, three name an id that no longer exists,
and the first real run died in 1.6 seconds with the provider's own message quoted.

\begin{lstlisting}[language=bash]
## Exit Criteria

- [ ] `cargo test -q -p pii-store --test retention_drops_what_the_policy_expires`
- [ ] `grep -q 'retention' crates/pii-store/src/retention.rs`
- [ ] `grep -q 'prune' crates/pii-cli/src/cmd_store.rs`
- [ ] `cargo test -q -p pii-store --test retention_keeps_a_live_run`
\end{lstlisting}

The criteria are commands, not descriptions, because three separate consumers parse them: the
judge runs them, the queue copies them into the row's verify column, and the gate's
\texttt{plan-verify} step checks that the copy still matches. A criterion written as prose is
a criterion no machine can hold a worker to.

\paragraph{The requirement a plausible implementation satisfies in appearance only.} Every
spec here carries one, and it is the most useful paragraph in the document. Spec 362's is that
a worker can add the config table, resolve the role in one node kind, and leave every other
call site hardcoded: the config exists, nothing uses it, and the next stale id hides. Writing
that paragraph forces the author to name the check that distinguishes the real implementation
from the plausible one.

\section{The row is the unit of dispatch}

A row is one line, seven tab-separated fields, no header:

\begin{longtable}{@{}p{0.10\textwidth}p{0.16\textwidth}p{0.66\textwidth}@{}}
\toprule
\textbf{Field} & \textbf{Name} & \textbf{What it is} \\
\midrule
\endhead
1 & \texttt{id} & the row's key in the state file; the spec number where one exists \\
2 & \texttt{spec} & the spec's file name under \texttt{specs/} \\
3 & \texttt{brief} & a path to the brief, or \texttt{-} \\
4 & \texttt{title} & one line, used in the commit message \\
5 & \texttt{verify} & the spec's exit criteria, joined with \texttt{;} \\
6 & \texttt{paths} & the paths this row may change, joined with \texttt{;} \\
7 & \texttt{gate} & the gate to run, or an explicit skip \\
\bottomrule
\end{longtable}

Column 5 is a \emph{cache} of column 2's document, not an independent copy. For an open row,
\texttt{state.py} re-derives it from the spec's \texttt{\#\# Exit Criteria} section every time
the plan is read; for a finished row the text is frozen, because regenerating it from a spec
amended later would claim the work was held to criteria that did not exist when it ran. The
reason is recorded in the code: 30 of 183 rows had dropped a criterion outright, and the drift
is invisible until the gate runs the wrong list.

Column 6 makes ``this work belongs to this row'' checkable. \texttt{row-paths.judge} compares
it with the spec's \texttt{\#\# Files} section, and \texttt{row-verify.judge} compares column
5 with the spec's criteria, command for command. Both refusals happen at pick time, before a
worker is bought.

\paragraph{Duplicate ids are refused, loudly.} The id is the key the state is stored under, so
two rows sharing one is not cosmetic: the second inherits the first's status and is never
picked. Measured once: two rows both said \texttt{A20}, and every cycle for about five hours
reported ``queue empty'' with a queued row sitting in the plan. The refusal names the
duplicates and says why.

\paragraph{A spec number must resolve to one file.} The dispatcher resolves
\texttt{specs/<N>-*.md} and the session recovery keys the resume handle by the same number.
Two files sharing a number send the row to the wrong document and resume another item's
finished session. Measured: a row was dispatched against a spec that was already committed,
resumed a ten-turn session, spent \$0.37, changed nothing, and then failed verify on its own
criterion.

\begin{center}
\resizebox{\linewidth}{!}{%
\begin{tikzpicture}[node distance=8mm and 12mm]
	\node[b] (spec) {spec\\\scriptsize premises $\cdot$ criteria $\cdot$ files};
	\node[b, right=of spec, yshift=12mm] (v) {column 5\\\scriptsize derived, open rows};
	\node[b, right=of spec, yshift=-12mm] (p) {column 6\\\scriptsize declared paths};
	\node[b, right=of v] (pick) {\texttt{pick}};
	\node[v, below=of pick] (match) {matches?};
	\node[b, right=of match] (row) {row dispatched};
	\node[b, left=of match, xshift=-2mm] (refuse) {refused\\\scriptsize named reason};

	\draw[f] (spec) -- (v);
	\draw[f] (spec) -- (p);
	\draw[f] (v) -- (pick);
	\draw[f] (pick) -- (match);
	\draw[f] (match) -- node[p, above=2pt, midway]{yes} (row);
	\draw[bk] (match) -- node[p, above=2pt, midway]{no} (refuse);
	\draw[sf] (p) -- (match);
\end{tikzpicture}}
\captionof{figure}{Two comparisons stand between the plan and a paid dispatch, and a row that
disagrees with its spec is refused with the field named while the state is left untouched.}
\end{center}

\section{The state file}

\texttt{scripts/autonomy/state.py} holds one JSON file, one object per row, keyed by id. Every
write goes to a temp file and is renamed into place, so a cycle killed mid-write cannot leave
a truncated state --- which is the whole reason it is not an append-only log.

Fields: \texttt{status}, \texttt{attempts}, \texttt{session}, \texttt{last},
\texttt{evidence}, and the two annotation fields \texttt{spec} and \texttt{commit}. Statuses
in use are \texttt{pending}, \texttt{in-progress}, \texttt{needs-review}, \texttt{done},
\texttt{skipped} and \texttt{parked}.

\paragraph{What \texttt{last} is for.} It is the reason a row is not offered again, written as
a short marker: \texttt{window-closed}, \texttt{worker-rc=3}, \texttt{gate-failed: <step>},
\texttt{verify-failed}, \texttt{precommit-broken}, \texttt{park-cause-repaired: <note>}. A row
that is parked or needs review waits for a person \emph{unless} its marker is one a session
can act on.

\begin{longtable}{@{}p{0.40\textwidth}p{0.16\textwidth}p{0.38\textwidth}@{}}
\toprule
\textbf{\texttt{last} marker} & \textbf{Resumable} & \textbf{Who acts next} \\
\midrule
\endhead
\texttt{window-closed} & yes, with the session & the next cycle inside the window \\
\texttt{worker-rc=N} & yes & the same session, resumed \\
\texttt{verify-failed} & yes & the same session, with a corrective brief \\
\texttt{gate-failed: <step>} & yes & the same session, with the failing step named \\
\texttt{precommit-broken}, \texttt{precommit-timeout} & yes & the same session; attempts not
spent \\
\texttt{park-cause-repaired: <note>} & yes & the repaired row goes back on the queue \\
anything else (e.g. \texttt{3-strikes}) & no & a person reads it \\
\bottomrule
\end{longtable}

\paragraph{Three strikes.} A row is parked when it reaches three attempts: at dispatch if the
count is already three, or at verify if the failure would make it three. The reason is written
as prose that states the rule --- the same mechanical failure repeated, so a fourth dispatch
would spend without new information. Attempts are never incremented by a timeout or a broken
pre-commit build, because neither is a verdict on the worker.

\paragraph{Repair, not reset.} \texttt{state.py repair <id> <note>} moves a parked row back to
\texttt{needs-review} (or \texttt{pending} without a session) and writes
\texttt{park-cause-repaired: <note>}. It deliberately does not touch attempts: the cap is a
safety property, so a repaired row that fails again parks again on its own merits.

\paragraph{Annotations are proven.} \texttt{state.py annotate <id> --spec <path> --commit
<sha>} refuses to write a spec that does not exist, a spec that is not the row's own, or a
commit git does not know. The field that can hold a lie is worse than the missing field: a row
naming a commit that does not contain its work is evidence an operator would act on.

\begin{lstlisting}[language=bash]
python3 scripts/autonomy/state.py pick | cut -f1          # the top dispatchable id, or NONE
python3 scripts/autonomy/state.py get 376                 # the row as JSON
python3 scripts/autonomy/state.py set 376 status=needs-review attempts=1 last=verify-failed
python3 scripts/autonomy/state.py annotate 376 --spec specs/376-retention-drops-what-the-policy-expires.md \
                                                --commit 0efc6fd
\end{lstlisting}

\section{The lane}

A dispatch goes through a named script, never an inline command. There are two lanes:

\begin{itemize}
	\item \textbf{coder} --- a one-node graph that runs a coder node on the spec. The default.
	\item \textbf{spec} --- the tracked graph \texttt{graphs/spec.toml}, which is the same
	implementation expressed as composition: \texttt{code} (a coder with its own wall and spend
	bounds), \texttt{verify} (a judge reading the spec at a commit), \texttt{docs} (a second
	coder whose brief is documentation only), \texttt{reverify} (the same criteria again, so a
	docs step that broke them is refused), then \texttt{done}, with \texttt{unverified} and
	\texttt{refused} as the two named refusals.
\end{itemize}

\texttt{PII\_LANE\_GRAPH=spec} selects the second; anything else keeps the first. The script
prints \texttt{LANE-GRAPH: <lane>} before anything is spent and again into the worker log, and
the cycle reads that line to record which lane a tree-changing run used. The commit carries it
as a \texttt{Lane:} trailer, so a commit says which of the two produced it. A run that changed
nothing leaves the previous record standing, because a commit must never be attributed to a
run that produced no tree.

\begin{center}
\resizebox{\linewidth}{!}{%
\begin{tikzpicture}[node distance=8mm and 12mm]
	\node[b] (row) {row};
	\node[b, right=of row] (disp) {dispatch script\\\scriptsize \texttt{LANE-GRAPH}};
	\node[b, right=of disp] (code) {\texttt{code}};
	\node[b, right=of code] (verify) {\texttt{verify}};
	\node[b, below=of verify] (docs) {\texttt{docs}};
	\node[b, left=of docs] (reverify) {\texttt{reverify}};
	\node[b, left=of reverify] (done) {\texttt{done}};
	\node[b, above=of reverify] (ref) {\texttt{unverified} $\cdot$ \texttt{refused}};

	\draw[f] (row) -- (disp);
	\draw[f] (disp) -- (code);
	\draw[f] (code) -- (verify);
	\draw[f] (verify) -- (docs);
	\draw[f] (docs) -- (reverify);
	\draw[f] (reverify) -- (done);
	\draw[bk] (verify) -- (ref);
	\draw[bk] (reverify) -- (ref);
\end{tikzpicture}}
\captionof{figure}{Two judges, one coder and one documentation pass in the spec lane, with
refusals wired by name so a failed implementation parks its row instead of ending the run with
the row still claimed.}
\end{center}

\section{One cycle, in order}

A cycle is deterministic on purpose: the reasoning belongs to the agent that calls it, and the
mechanics belong in a script that can be tested. \texttt{scripts/autonomy/cycle.sh} does the
following, in this order, and the order is load-bearing:

\begin{enumerate}
	\item \textbf{The tree hold and the lock.} A background cycle yields if an interactive turn
	holds the tree, and reports \texttt{OUTCOME=tree-held}. The lock spans the whole cycle;
	\texttt{-o} is passed to \texttt{flock} so a daemonized grandchild cannot inherit the
	descriptor and hold the lock after the cycle is gone.
	\item \textbf{The operator's switch.} \texttt{PAUSED} in the control file stops before
	spending.
	\item \textbf{The window probe.} A closed account window exits 3 with nothing spent and the
	row not advanced. An undecided probe exits 4 rather than guessing.
	\item \textbf{The pick.} \texttt{state.py pick} returns the top dispatchable row, or
	\texttt{NONE}, or a refusal. A plan that cannot be read is a loud failure, never an empty
	queue: ``nothing pending'' and ``the queue is broken'' are different outcomes and only one
	is safe to sleep on.
	\item \textbf{The re-read.} Between the pick and the dispatch, the state is read again. A
	row that another cycle closed is not dispatched; a row whose marker says the work is
	already \texttt{worker-ok}, \texttt{verified} or \texttt{gated} resumes at verify with no
	worker bought.
	\item \textbf{The tree.} A dirty tree is handled in three moves: save, attribute, clean. A
	tree the loop's own last row left is stashed so this gate measures only this change; a tree
	no row owns is a person's, and the cycle stops with \texttt{OUTCOME=dirty-tree}.
	\item \textbf{The dispatch.} A session id recovered from the worker's own log is used as a
	resume handle if one exists; otherwise the worker is dispatched with the lane script, and
	the session id is recorded.
	\item \textbf{The verify.} The row's own commands, run in a login shell that sources the
	loop's environment first. Measured: a login shell resets \texttt{PATH} for root and loses
	\texttt{/usr/local/cargo/bin}, so two cargo criteria were recorded as failures on a tree
	where all four passed.
	\item \textbf{The gate.} Only an explicit \texttt{0}, \texttt{-} or \texttt{none} skips it.
	A green verify is not a green gate.
	\item \textbf{The pre-commit build.} The door builds \texttt{HEAD} plus the index before
	the commit exists, so a tree that does not build is never committed.
	\item \textbf{The commit, then the head build.} The commit is made with an explicit neutral
	identity; then \texttt{head-builds.sh} builds the commit that was just made, because the
	gate measured the working tree and this reads the commit.
	\item \textbf{The annotation and the chain.} The row records its spec and its commit; the
	next row starts as soon as the lock is free.
\end{enumerate}

\begin{center}
\resizebox{\linewidth}{!}{%
\begin{tikzpicture}[node distance=8mm and 12mm]
	\node[b] (pending) {pending};
	\node[b, right=of pending] (running) {in-progress};
	\node[b, right=of running, yshift=12mm] (review) {needs-review};
	\node[b, right=of running, yshift=-12mm] (done) {done};
	\node[b, right=of review] (parked) {parked};
	\node[v, below=of review] (cap) {cap?};

	\draw[f] (pending) -- node[p, above=2pt, midway]{pick} (running);
	\draw[f] (running) -- node[p, above=2pt, midway]{criteria + gate} (done);
	\draw[f] (running) -- node[p, above=2pt, midway]{failed} (review);
	\draw[f] (review) -- node[p, right=2pt, midway]{attempts} (cap);
	\draw[f] (cap.west) -| node[p, below=2pt, midway]{resume} (running.south);
	\draw[bk] (review) -- node[p, above=2pt, midway]{cap reached} (parked);
	\draw[sf] (parked.south) |- node[p, below=2pt, midway]{repair} (pending.south);
\end{tikzpicture}}
\captionof{figure}{Four statuses and one decision: a row that fails goes to needs review with
a marker, resumes in the same session while attempts remain, and parks for a person at three.}
\end{center}

\section{Resume before retry}

The rule is old and it is the cheapest in the system: continue the existing session. A fresh
start re-derives context the previous run already established and already billed for.

Mechanically, a resume needs a session id, and the session id is unrecoverable once the
process exits. So it is written at dispatch time, and the recovery path is layered:

\begin{itemize}
	\item The dispatch script records the id in \texttt{\_tmp/claude-batchC/session-<N>-*.txt}.
	\item If the process was cut off before that, \texttt{session.sh recover <N>} reads the
	worker's own log, and the cycle records the id with \texttt{last=resumed-from-worker-log}.
	\item A worker cut off by a closed account window records its session before exiting, so
	the next cycle resumes rather than restarting from scratch. Without that, a dirty tree with
	no session was kept as a leftover and the spec was dispatched again from zero.
\end{itemize}

\paragraph{A resume carries a diagnosis.} When the previous attempt ended
\texttt{gate-failed}, \texttt{verify-failed} or \texttt{worker-rc}, the resume is sent with a
corrective brief that names what failed and includes the failing output. The worker cannot see
the gate's verdict, and its own criteria may well have passed --- which is exactly how one
first pass went: criteria 4 of 4, gate red on a tree the export had not included. Naming the
failure is the difference between a fix and a retry.

The brief's evidence block has its own history. It was built from the newest cycle log, which
was usually the accumulator file being appended to while the brief was assembled --- so a
brief for one row quoted rows from a previous month. It now takes the newest log that is
neither this cycle's nor the accumulator, and falls back to the log's tail so it can never
print nothing.

\paragraph{What is not a row fault.} Three failure classes are separated from the row on
purpose:

\begin{itemize}
	\item \textbf{Quota.} A provider limit puts the row back on the queue with its attempts and
	session intact, and the reset time the CLI printed is recorded. Never a blind retry, never
	a scope cut.
	\item \textbf{Authentication.} A dead login refuses every row, so counting it as an attempt
	would park three rows that were never tried. The row keeps its attempts and the queue is
	paused for a person.
	\item \textbf{A killed gate.} A timeout exits with its own outcome; attempts are not
	incremented and the marker keeps the next cycle at verify, so a slow host costs a retry and
	never a re-dispatch.
\end{itemize}

\begin{center}
\resizebox{\linewidth}{!}{%
\begin{tikzpicture}[node distance=8mm and 12mm]
	\node[v] (why) {why?};
	\node[b, right=of why, yshift=14mm] (win) {window closed\\\scriptsize resume in window};
	\node[b, right=of why, yshift=4mm] (quota) {quota\\\scriptsize requeue, same session};
	\node[b, right=of why, yshift=-6mm] (auth) {auth\\\scriptsize pause the queue};
	\node[b, right=of why, yshift=-16mm] (fail) {gate or verify failed\\\scriptsize corrective brief};
	\node[b, below=of why] (fresh) {no session anywhere\\\scriptsize a fresh dispatch};

	\draw[f] (why) -- (win);
	\draw[f] (why) -- (quota);
	\draw[f] (why) -- (auth);
	\draw[f] (why) -- (fail);
	\draw[sf] (why) -- node[p, right=2pt, midway]{never recorded} (fresh);
\end{tikzpicture}}
\captionof{figure}{Four resumable reasons and one that is not, and only the last one justifies
paying for context twice.}
\end{center}

\section{Quarantine: a park must not wedge the queue}

Parking used to leave the row's working-tree changes in place, and the dispatcher refuses to
dispatch on a dirty tree, so one parked row stopped every later row for a whole night.
Measured: a row parked with nine files and the next row never ran --- and the parked work was
good, five tests green.

\texttt{scripts/autonomy/quarantine-wip.sh <row> <reason>} now puts the changes aside without
deleting anything:

\begin{enumerate}
	\item Every changed path is \emph{copied} to
	\texttt{/projects/\_trash/<row>-<reason>-<ts>/} with the tree layout preserved, and
	\texttt{git status --porcelain} is saved as \texttt{status.txt}.
	\item A modified tracked file is restored to its \texttt{HEAD} content with \texttt{git
	show HEAD:<path>} --- explicitly \emph{not} \texttt{git checkout --}, which is banned in
	this workspace because it destroys uncommitted work.
	\item An untracked file is \emph{moved} out of the tree, never removed.
	\item The directory is recorded in \texttt{ASIDE.txt}, and the path is printed on stdout so
	the cycle can put it in the row's \texttt{last} note.
\end{enumerate}

If anything is still dirty afterwards the script exits 3 with \texttt{STILL-DIRTY} and the
count, and the cycle says so: the dispatcher will refuse until the tree is clean, and a
warning is better than a silent night.

\begin{center}
\resizebox{\linewidth}{!}{%
\begin{tikzpicture}[node distance=8mm and 12mm]
	\node[b] (dirty) {dirty tree};
	\node[b, right=of dirty] (copy) {copy every path\\\scriptsize aside, layout kept};
	\node[b, right=of copy, yshift=10mm] (mod) {restore from HEAD\\\scriptsize \texttt{git show}};
	\node[b, right=of copy, yshift=-10mm] (unt) {move untracked\\\scriptsize \texttt{mv}};
	\node[v, below=of copy] (clean) {clean?};
	\node[b, left=of clean] (ok) {queue moves\\\scriptsize row parked};
	\node[b, right=of clean] (still) {STILL-DIRTY\\\scriptsize exit 3};

	\draw[f] (dirty) -- (copy);
	\draw[f] (copy) -- (mod);
	\draw[f] (copy) -- (unt);
	\draw[f] (mod) -- (clean);
	\draw[f] (unt) -- (clean);
	\draw[f] (clean) -- node[p, above=2pt, midway]{yes} (ok);
	\draw[bk] (clean) -- node[p, above=2pt, midway]{no} (still);
\end{tikzpicture}}
\captionof{figure}{Copy first and restore second, because a restore that loses the work it was
preserving is the one unrecoverable path.}
\end{center}

\begin{lstlisting}[language=bash]
$ bash scripts/autonomy/quarantine-wip.sh 352 parked
QUARANTINED 9
/projects/_trash/352-parked-<UTC timestamp>
$ git status --porcelain | wc -l
0
\end{lstlisting}

\section{Review rows}

A review is work with a brief, a deliverable and a verifier, so it is a row. The user
directive behind it was that an analysis must also be put in the queue, positioned before the
items that matter.

A review spec carries two extra lines, \texttt{BRIEF: <path>} and \texttt{DELIVERABLE:
<path>}. \texttt{review-row.sh} reads both, refuses if the brief is missing or the deliverable
line is absent, dispatches the read-only reviewer in plan mode, then copies the newest
\texttt{.md} in the brief's directory to the deliverable and exits 5 if the deliverable is
missing. The check is on the artefact, not on the reviewer's claim to have produced one.

\section{Commits}

A commit is one message saying what the change does and why the alternative was rejected.
Three rules apply to every commit made by the loop:

\begin{enumerate}
	\item \textbf{An explicit neutral identity.} The repository's own git configuration may
	carry a tool name, and then every plain \texttt{git commit} breaches the rule silently. The
	cycle passes \texttt{-c user.name} and \texttt{-c user.email} on the command line, so the
	identity is a property of the commit and not of the machine.
	\item \textbf{Only the row's paths.} \texttt{git add} is given the row's declared paths,
	and \texttt{unlisted-changes.sh} refuses the commit if the index would include a file the
	row does not list. The tree is left as the worker made it, because the next cycle resumes
	at verify.
	\item \textbf{The gate's verdict is in the message, and the lane is a trailer.} The message
	states that the criteria are in the spec and the gate is green, and a \texttt{Lane:
	spec|coder} trailer names the lane when one was recorded.
\end{enumerate}

\paragraph{Nothing staged is not nothing to do.} If the add leaves the index empty, the work
may already have landed in an earlier commit. The cycle closes the row on exactly the evidence
the normal path uses --- the row's own criterion plus a green gate --- and names the commit
that carries the work. A leftover commit is still a commit; re-dispatching it spends the same
spec twice.

\begin{lstlisting}[language=bash]
git add -- crates/pii-store/src/retention.rs crates/pii-store/tests/...
UNLISTED=$(bash scripts/autonomy/unlisted-changes.sh "$PATHS") || { git reset -q; exit 17; }
git -c user.name="Marcelo Correa" -c user.email="mvcorrea@gmail.com" commit -q \
  -m "376-retention: automated pass - retention drops what the policy expires (gate green)

Lane: spec"
\end{lstlisting}

The identity is on the command line rather than in the repository's configuration for the
reason recorded in the workspace rules: a repository whose own \texttt{git config} holds a
tool name breaches the rule on every plain commit, silently, because nobody re-reads the
config. That happened for thirty-five commits in one repository here before anyone looked.

\section{The review that gave the backlog verdicts}

The parked backlog was reviewed on 2026-11-04 under spec 377, and the method is worth copying:
every criterion of every parked row was \emph{executed}, from the plan's column 5, split on
\texttt{;} the same way the runner splits it. Nothing was read from a helper's summary.

The review produced four kinds of verdict, and the distinctions are the useful part:

\begin{itemize}
	\item \texttt{proceed} --- the row's work is not on disk, so it is a clean dispatch from
	\texttt{HEAD}.
	\item \texttt{close} --- the row was superseded by later work, and its criteria pass or are
	stale. Closing it is a state write, not code.
	\item \texttt{needs-decision} --- the row cannot be judged without an answer only the
	operator can give (a budget question, or which of two commits answers the row).
	\item missing evidence --- a WIP directory the spec names does not exist, so any verdict
	depending on it is unverified rather than done.
\end{itemize}

It also named four failure classes, each of which recurs in this manual: a guard that fires on
the case it was built to permit; criteria that name files which no longer exist; a sandbox in
which the criteria cannot run as written (a read-only \texttt{/tmp}, no \texttt{uv} cache, no
\texttt{python3}); and quarantined work that is not where the spec says it is.

\section{Failure modes worth recognising}

\begin{enumerate}
	\item \textbf{A queue that reports itself empty while work is queued.} Cause: two rows
	sharing an id, so the second inherits the first's status. The plan refuses duplicates now.
	\item \textbf{A night with one row parked and nothing else done.} Cause: a park that left
	the tree dirty. Quarantine is the fix, and the row's \texttt{last} note carries the aside
	path.
	\item \textbf{The same spec bought twice.} Cause: a resume that did not happen, or a number
	that resolved to the wrong file. Check the session file and the spec number before the next
	dispatch.
	\item \textbf{A criterion that passes on the worker's tree and fails in the gate.} Cause:
	the export did not include the change (a deletion, or a file the overlay did not see). The
	gate's export-lost and overlay-deletion handling exist for this.
	\item \textbf{A commit that names a tool.} Cause: the repository's git configuration. Read
	it, and pass the identity explicitly.
	\item \textbf{A worker that fixes the symptom.} Cause: a corrective brief that named the
	failing command but not the rule. The brief says: fix the cause, do not weaken or delete an
	assertion, do not raise a timeout, do not add a sleep, do not touch anything outside the
	spec's scope.
\end{enumerate}
